\chapter{The Imperfect --- the Past That Kept Going}
\label{ch:imperfect}

One past tense says a thing finished; this one says a thing was
\emph{going on}. It is also the most regular tense in the language, with exactly
three irregular verbs in the whole of Spanish --- and then one genuinely new
distinction to learn to think in.
\emph{Practice lessons: \texttt{lessons/ES-C16-*.md}.}

\section{\emph{hablaba}, \emph{com\'{i}a} --- the past that kept going}
\label{lesson:imperfecto}

A \textbf{tense} is the shape a verb takes to place an action in time. Spanish
has two ordinary past tenses. The \textbf{preterite} reports that a past thing
\textbf{finished}: \emph{habl\'{e}}, ``I spoke.'' The \textbf{imperfect} reports
that a past thing was \textbf{going on}, or happened over and over: ``I
\textbf{was speaking},'' ``I \textbf{used to speak}.'' That is what this whole
chapter is about.

And here is the good news: the imperfect is the \textbf{most regular tense in
the language}. Drop the infinitive ending (\emph{-ar}, \emph{-er}, \emph{-ir})
and add one of two sets of endings.

\begin{center}
\begin{tabular}{lll}
\toprule
Person & \emph{-ar} (\emph{hablar}) & \emph{-er} / \emph{-ir} (\emph{comer}, \emph{vivir}) \\
\midrule
yo & habl\textbf{aba} & com\textbf{\'{i}a} / viv\textbf{\'{i}a} \\
t\'{u} & habl\textbf{abas} & com\textbf{\'{i}as} \\
\'{e}l / ella / usted & habl\textbf{aba} & com\textbf{\'{i}a} \\
nosotros & habl\textbf{\'{a}bamos} & com\textbf{\'{i}amos} \\
ellos / ustedes & habl\textbf{aban} & com\textbf{\'{i}an} \\
\bottomrule
\end{tabular}
\end{center}

The \textbf{\emph{-er} and \emph{-ir} verbs share one set again} --- the third
time now (they also shared in the preterite). Two conjugations, one pattern.

\begin{cousinweb}
\textbf{-aba} $\leftarrow$ Latin \emph{-\={a}bam} $\cdot$ \textbf{-\'{i}a}
$\leftarrow$ Latin \emph{-\={e}bam}.\\[3pt]
\textbf{Same Latin ending, two fates.} The \emph{b} survived whole in
\emph{-aba} (\emph{-\={a}bam} $\to$ \emph{-aba}), but wore away in the other set
(\emph{-\={e}bam} $\to$ \emph{-ea} $\to$ \emph{-\'{i}a}). One sound, two
outcomes, inside the same language.
\end{cousinweb}

\begin{sounds}
\textbf{accent-acute} --- \emph{habl\'{a}bamos}: the stress falls
\textbf{three syllables back} (\emph{ha-BL\'{A}-ba-mos}), and Spanish writes an
accent whenever it does. This is the \emph{only} form in the \emph{-ar} set that
takes one.\\[3pt]
\textbf{hiatus-ia} --- \emph{com\'{i}a}: here the accent \textbf{splits the
vowels}. Without it, \emph{ia} would be a single syllable (as in \emph{piano});
with it, \emph{co-m\'{i}-a} is three. \textbf{Every} form in the
\emph{-er}/\emph{-ir} set needs it.
\end{sounds}

\begin{grammarlens}
\textbf{One trap, and a broken rule.} \emph{Yo hablaba} and \emph{\'{e}l
hablaba} are \textbf{identical} --- the ending no longer tells you who is
speaking. Spanish normally \emph{drops} its subject pronouns, because the ending
carries the person; in the imperfect the ending stopped doing that job, so the
pronoun steps back in. You will hear \emph{yo} said out loud here far more often
than in the present tense.
\end{grammarlens}

Say them straight through: \emph{hablaba, hablabas, hablaba, habl\'{a}bamos,
hablaban} $\cdot$ \emph{com\'{i}a, com\'{i}as, com\'{i}a, com\'{i}amos,
com\'{i}an}. And in a sentence: \emph{Cuando era ni\~{n}o, hablaba espa\~{n}ol
en casa} (``When I was a child, I used to speak Spanish at home'').

\section{\emph{era}, \emph{iba}, \emph{ve\'{i}a} --- all three of them}
\label{lesson:tres-irregulares}

The preterite hands you a pile of strong, irregular forms to memorise. Here is
the compensation: the imperfect has \textbf{three} irregular verbs. Not three
families --- \textbf{three verbs}, in the entire language.

\begin{center}
\begin{tabular}{llll}
\toprule
Person & \emph{ser} $\to$ \textbf{era} & \emph{ir} $\to$ \textbf{iba} & \emph{ver} $\to$ \textbf{ve\'{i}a} \\
\midrule
yo & era & iba & ve\'{i}a \\
t\'{u} & eras & ibas & ve\'{i}as \\
\'{e}l / ella / usted & era & iba & ve\'{i}a \\
nosotros & \'{e}ramos & \'{i}bamos & ve\'{i}amos \\
ellos / ustedes & eran & iban & ve\'{i}an \\
\bottomrule
\end{tabular}
\end{center}

That is the entire irregular inventory. Every other verb in Spanish ---
thousands of them --- takes the regular endings from the last section.

And \emph{ver} barely counts: it simply keeps an extra \emph{-e-} from its older
form \emph{veer}, then conjugates normally (\emph{ve-} $+$ \emph{-\'{i}a}).

\begin{cousinweb}
\textbf{era} $\leftarrow$ Latin \emph{eram}, from the same \emph{es-} root that
gives \emph{ser}, \emph{soy}, \emph{es} --- and English \textbf{is} and
\textbf{are}. So \emph{ser} uses that ancient root in the present \emph{and} the
imperfect, but jumps to a different verb entirely, Latin \emph{fu\={\i}}, in the
preterite: \emph{era} versus \emph{fue}.
\end{cousinweb}

\begin{cousinweb}
\textbf{\emph{ir}'s third Latin verb.} Spanish \emph{ir} (``to go'') is
stitched together out of \textbf{three different Latin verbs}, one per
tense:\\[3pt]
present \emph{voy, vas, va} $\leftarrow$ Latin \emph{v\={a}dere}, ``to
advance''\\[3pt]
preterite \emph{fui, fuiste, fue} $\leftarrow$ Latin \emph{fu\={\i}}, from
\emph{esse}\\[3pt]
imperfect \textbf{\emph{iba, ibas, iba}} $\leftarrow$ Latin \emph{\={\i}re} ---
\emph{ir}'s \textbf{own infinitive}\\[3pt]
The imperfect is the \textbf{one place} the original \emph{\={\i}re} survives
conjugated. The infinitive \emph{ir} and the imperfect \emph{iba} are the last
two pieces of the real verb; everything else was borrowed from elsewhere.
\end{cousinweb}

\begin{sounds}
\textbf{stress-antepenultimate} --- \textbf{\'{e}ramos} and
\textbf{\'{i}bamos} take a written accent for exactly the reason
\emph{habl\'{a}bamos} does: the stress lands three syllables back
(\emph{\'{E}-ra-mos}, \emph{\'{I}-ba-mos}). Those two are the only accented
forms in the irregular set apart from \emph{ver}, whose accent is the
vowel-splitting kind (\emph{ve-\'{i}-a}).
\end{sounds}

Say them: \emph{era, eras, era, \'{e}ramos, eran} $\cdot$ \emph{iba, ibas, iba,
\'{i}bamos, iban}. In a sentence: \emph{Cuando era ni\~{n}o, iba a la escuela}
(``When I was a child, I used to go to school''). And say the three sources out
loud: \emph{voy} $\leftarrow$ \emph{v\={a}dere} $\cdot$ \emph{fui} $\leftarrow$
\emph{fu\={\i}} $\cdot$ \emph{iba} $\leftarrow$ \emph{\={\i}re}.

\section{Practice --- two pasts, one choice}
\label{lesson:ch16-practice}

You now have \textbf{both} past tenses. English marks this difference only
sometimes, and never with an ending --- so this is a genuinely \textbf{new
distinction} to think in, not merely new forms to learn.

\begin{center}
\begin{tabular}{lll}
\toprule
 & preterite & imperfect \\
\midrule
what it does & a \textbf{completed} event & an \textbf{ongoing} or \textbf{habitual} state \\
shape of time & a \textbf{point} & a \textbf{line} \\
English & ``I ate'' & ``I \textbf{was} eating'' / ``I \textbf{used to} eat'' \\
typical clues & \emph{ayer, una vez, de repente} & \emph{siempre, todos los d\'{i}as, mientras} \\
\bottomrule
\end{tabular}
\end{center}

\textbf{Com\'{i}} $=$ the meal happened and ended. \textbf{Com\'{i}a} $=$ I was
in the middle of it, or I did it regularly.

\begin{grammarlens}
\textbf{Background and interruption.} The classic pairing: the imperfect sets
the scene, the preterite breaks it.\\[3pt]
\emph{\textbf{Hablaba} con Mar\'{i}a cuando \textbf{lleg\'{o}} Juan.} (``I was
talking to Mar\'{i}a when Juan arrived.'')\\[3pt]
The talking is the \textbf{line}; the arriving is the \textbf{point} that lands
on it. Swap the two tenses and the sentence stops making sense.
\end{grammarlens}

\begin{culture}
\textbf{Verbs that change meaning.} With a few verbs the choice does not just
shift the shape of time --- it changes \textbf{what the word means}. The logic
is consistent: the imperfect gives you the \emph{state}, the preterite gives you
the \emph{instant the state began}. Knowing $\to$ the moment of finding out.
\end{culture}

\begin{center}
\begin{tabular}{lll}
\toprule
verb & imperfect (state) & preterite (the moment it started) \\
\midrule
\emph{saber} & \emph{sab\'{i}a} --- I \textbf{knew} & \emph{supe} --- I \textbf{found out} \\
\emph{conocer} & \emph{conoc\'{i}a} --- I \textbf{knew} (a person) & \emph{conoc\'{i}} --- I \textbf{met} \\
\emph{tener} & \emph{ten\'{i}a} --- I \textbf{had} & \emph{tuve} --- I \textbf{got, received} \\
\emph{querer} & \emph{quer\'{i}a} --- I \textbf{wanted} & \emph{quise} --- I \textbf{tried} \\
\bottomrule
\end{tabular}
\end{center}

(And the negative sharpens it further: \emph{no quise} $=$ ``I
\textbf{refused}.'')

\textbf{Drill --- which tense?} Decide before you read the answer.

\begin{center}
\begin{tabular}{ll}
\toprule
sentence & tense \\
\midrule
``Every summer we \textbf{went} to Spain.'' & imperfect (\emph{\'{i}bamos}) --- habitual \\
``Yesterday we \textbf{went} to Spain.'' & preterite (\emph{fuimos}) --- one event \\
``She \textbf{was} tall.'' & imperfect (\emph{era}) --- a lasting state \\
``The concert \textbf{was} yesterday.'' & preterite (\emph{fue}) --- a bounded event \\
\bottomrule
\end{tabular}
\end{center}

Say these aloud: \emph{Com\'{i}a cuando llamaste} (``I was eating when you
called'') $\cdot$ the pair \emph{sab\'{i}a} ``I knew'' \ldots\ \emph{supe} ``I
found out'' $\cdot$ \emph{Cuando era ni\~{n}o, iba a la playa todos los
veranos}.

Recall it all: preterite or imperfect for a \textbf{habit}?
(\textbf{Imperfect}.) Which tense sets the background, and which interrupts it?
(Imperfect \textbf{sets}, preterite \textbf{interrupts}.) What is the difference
between \emph{conoc\'{i}a} and \emph{conoc\'{i}}? (``I \textbf{knew} someone''
versus ``I \textbf{met} them.'') How many irregular imperfects are there?
(\textbf{Three}.) And the one-word test, the thing to keep:
\textbf{line} $\to$ imperfect; \textbf{point} $\to$ preterite. Next chapter:
talking about what \textbf{will} happen.
